\begin{frame}{13.3 자주 하는 실수}
  \begin{enumerate}
    \item \textbf{학습 중 노이즈 있는 에피소드 리턴을 정책 성능으로
      판단}. 두 곡선 현상(확률적 악화, 결정론적 개선)은 오진의
      전형적 원천 --- ``노이즈 에피소드''만 보면 ``학습할수록
      나빠졌다''고 결론 내고 \emph{구 weights}로 돌아가는 실수.
      규칙: \textbf{탐색용 학습 곡선 $\neq$ 배포 품질 곡선}.
    \item \textbf{PPO가 손제어기와 비슷해지면 ``구현 버그''라고 의심}.
      점검 순서: (1) \emph{예산} --- 100만 스텝은 장난감, 논문은
      수백만$\sim$수천만. (2) \emph{구조적 우위} --- 1줄 제어기는
      오차 벡터 obs[8:10]을 \emph{공짜로} 받아 쓰는데 PPO는 시행착오로
      \emph{배워야}. (3) $\sigma$ 클램프. 실전 표준은 손제어기를
      \emph{베이스라인}으로 두고 RL이 \emph{추가}로 이득을 주는지
      재는 것.
    \item \textbf{$\sigma$ 상한을 두지 않고 내버려 두기}. 클램프 없으면
      엔트로피 보상이 $\sigma$ 를 2, 3으로 밀어올리고, \textbf{샘플의 상당
      부분이 범위 밖으로 나가 \texttt{np.clip}에게 잘림} --- 이 순간
      실제로 실행되는 분포는 가우시안이 아니라 \emph{잘려 변형된}
      분포. 클립은 ``거의 발생하지 않아야'' 정상 신호.
  \end{enumerate}
\end{frame}
